\section{Formal problem statement and method equations}
\label{app:formal}

This appendix states the benchmark and the reference methods in symbols. It
adds no new claim. Every quantity is the one the released code computes, and
Table~\ref{tab:app-settings} lists the settings.

\subsection{Design space, evaluator, and score}

A design is a vector $z \in \mathcal{B} = [-3.5, 3.5]^{38}$ of standardized
coordinates. Its physical parameters are $x = \mu + \sigma \odot z$, where
$\mu$ and $\sigma$ hold the per-coordinate mean and standard deviation of the
SuperWing designs and $\odot$ is the coordinatewise product. A task is a target
lift coefficient $C_L^{*} \in \{0.15, 0.41, 0.66\}$. Every design is evaluated
at three flow conditions $c \in \{1, 2, 3\}$ with Mach numbers $0.7853$,
$0.8195$, and $0.8530$ and Reynolds number $2 \times 10^{7}$, as in
Section~\ref{sec:benchmark}.

For one condition the evaluator meshes the wing, adjusts the angle of attack
$\alpha$ until the computed lift coefficient matches the target, and returns
the lift-to-drag ratio at that angle,
\begin{equation}
r_c(x) = \frac{C_L(x, c, \alpha_c^{*})}{C_D(x, c, \alpha_c^{*})},
\qquad
\bigl|C_L(x, c, \alpha_c^{*}) - C_L^{*}\bigr| \le 0.02 .
\label{eq:condition}
\end{equation}
The condition is valid only if a mesh exists, the solver process exits
normally, the solver's lift oscillation measure is below $5 \times 10^{-4}$, the
trimmed angle lies in $[-10^{\circ}, 25^{\circ}]$, and $0 < r_c(x) \le 40$.
Otherwise the condition returns a failure type instead of a number. The score
is the geometric mean over the three conditions,
\begin{equation}
y(x) = \Bigl( \prod_{c=1}^{3} r_c(x) \Bigr)^{1/3},
\label{eq:score}
\end{equation}
and it exists only when all three conditions are valid. We write
$d(x) \in \{0, 1\}$ for this event, delivery. When $d(x) = 0$ the
evaluation returns one of three outcome types. A conclusive failure means the
geometry was rejected or a condition was conclusively invalid. Partial means
some conditions returned values and none was invalid. Undetermined means no
condition returned one. In the released records every conclusive failure is a
rejected geometry.

The surrogate $s(x)$ is the same geometric mean computed from AeroTransformer
predictions with the same trim procedure. One surrogate evaluation takes about
a second and returns a number for every finite $x$, including designs with
$d(x) = 0$.

\subsection{An episode}

An episode is one search on one task. Let $\mathcal{H}_k$ denote everything the
optimizer has observed after $k$ online evaluations, that is, its surrogate
queries, the designs it submitted, and for each submitted design either the
score or the failure type. A policy $\pi$ maps $\mathcal{H}_k$ either to a
batch of at most 16 surrogate queries or to the next online design $x_{k+1}$.
The budgets are
\begin{equation}
\text{surrogate evaluations} \le 4096, \qquad k \le 32,
\label{eq:budgets}
\end{equation}
and online evaluations are strictly sequential, so $x_{k+1}$ may depend on
$d(x_k)$ and $y(x_k)$. At each checkpoint $q \in \mathcal{Q} = \{0, 1, 2, 4, 8,
16, 32\}$ the policy also names one recommendation $\rho_q \in \mathcal{B}$,
the design it would hand to a user. The recommendation is graded by the same
evaluator, the grade is never shown to the policy, and it costs no budget. If
$\rho_q$ was already evaluated online, the grade reuses that result. A design
submitted twice in one episode returns its recorded result and still costs a
query.

\subsection{Endpoints}

For one method and one task with episodes $e = 1, \dots, n$, where $n = 10$
seeds, write $d_q^{e} = d(\rho_q^{e})$ and $y_q^{e} = y(\rho_q^{e})$. Then
\begin{equation}
D_q = \frac{1}{n} \sum_{e=1}^{n} d_q^{e},
\qquad
Q_q = \frac{\sum_{e} d_q^{e}\, y_q^{e}}{\sum_{e} d_q^{e}},
\qquad
O = \frac{1}{32\, n} \sum_{e=1}^{n} \sum_{k=1}^{32} d(x_k^{e}).
\label{eq:endpoints}
\end{equation}
$D_q$ is recommendation delivery, $Q_q$ is conditional quality, defined only
when some recommendation delivers, and $O$ is online delivery. The paper
reports $D_q$ and $Q_q$ at $q = 32$ for the baselines and at $q = 4$ for the
agents. The delivery averaged across tasks in Section~\ref{sec:analysis} is the
mean of $D_{32}$ over tasks.

\subsection{The belief model of the surrogate-assisted baselines}

The four surrogate-assisted baselines share one belief model, fitted anew in
every episode from the 62 warm-start results and the episode's own online
outcomes. Its feature map is
\begin{equation}
\phi(x) = \bigl( 1,\; z_1/3.5,\; \dots,\; z_{38}/3.5,\; \log s(x) - \log 30
\bigr) \in \mathbb{R}^{40}.
\label{eq:features}
\end{equation}
The delivery model is a weighted logistic regression with an $\ell_2$ penalty,
\begin{equation}
\begin{aligned}
p(x) &= \operatorname{clip}\bigl( \sigma(w^{\top} \phi(x)),\, 0.05,\, 0.95 \bigr), \\
w &= \arg\min_{w} \sum_{i} \omega_i\, \ell\bigl( d_i, \sigma(w^{\top} \phi(x_i)) \bigr)
+ 4\, \| w_{1:39} \|^{2} + 0.1\, w_0^{2},
\end{aligned}
\label{eq:delivery-model}
\end{equation}
where $\ell$ is the logistic loss, the label $d_i$ is 1 for a pass and 0
otherwise, and the weight $\omega_i$ is 1 for a pass or a conclusive failure
and $0.5$ for a partial or undetermined outcome. Appendix~\ref{app:weight}
varies this weight. The fit runs 40 iterations. The discrepancy model is a
ridge regression, on the delivered results, of the log ratio between the CFD
score and the surrogate score,
\begin{equation}
\begin{aligned}
\delta_i &= \log y_i - \log s(x_i), \qquad \hat{\delta}(x) = v^{\top} \phi(x), \\
v &= \arg\min_{v} \sum_{i:\, d_i = 1} \bigl( \delta_i - v^{\top} \phi(x_i) \bigr)^{2}
+ 10\, \| v_{1:39} \|^{2} + 0.1\, v_0^{2},
\end{aligned}
\label{eq:discrepancy-model}
\end{equation}
with a predictive standard deviation $\hat{\sigma}(x)$ floored at $0.05$ and
set to $0.25$ while no delivered result exists. The private merit of a design
combines the two models,
\begin{equation}
m(x) = \log p(x) + \log s(x) + \hat{\delta}(x) + \tfrac{1}{2}\, \hat{\sigma}(x)
+ \tfrac{1}{10}\, H\bigl( p(x) \bigr),
\qquad
H(p) = -p \log p - (1 - p) \log (1 - p).
\label{eq:merit}
\end{equation}
The first three terms are the log of the expected delivered score, the
delivery probability times the corrected surrogate score. The last two terms
reward uncertainty in value and delivery.

\subsection{The five baseline variants}

All four surrogate-assisted variants run their generator in unit coordinates
on the box with a population of 16 and 256 generations, which is the
4{,}096-evaluation budget, and keep an active archive of the 64 best
cheap-valid designs. A design whose surrogate evaluation is not finite receives
a loss of $10^{12}$.

\textbf{CMA-ES, no feedback.} The objective is $s(x)$. All 256 generations run
before the first online query. The 32 best distinct archive designs are
queued in descending $s$, and online results change nothing.

\textbf{CMA-ES, selection only.} The same search and the same archive. At online
step $k$ the method submits $\arg\max m(x)$ over the unqueried archive designs,
refitting $p$ and $\hat{\delta}$ after every result. The archive never
changes, so feedback picks the next design but cannot create a new one.

\textbf{Feedback-guided CMA-ES.} One continuing CMA-ES state with initial mean
$0.5$ in unit coordinates and initial step size $\sigma_0 = 0.25$. Before each
online query it runs eight generations, 128 surrogate evaluations, with
objective $m(x)$, then submits the archive $\arg\max$ of $m$. Both models
are refitted after every result, so feedback changes the selection and the
proposal distribution.

\textbf{Feedback-guided DE.} DE/rand/1/bin with $F = 0.8$ and
$\mathit{CR} = 0.9$ replaces the CMA-ES state. A
trial replaces its parent when its merit is at least the parent's, and a
mutant coordinate outside the box is redrawn uniformly. Schedule and selector
are those of feedback-guided CMA-ES.

\textbf{SMAC3.} The black-box facade of SMAC 2.4.0 with its Gaussian-process
model, 8 Sobol initial designs, 24 model-guided queries, one evaluation per
configuration, and a common penalty value for every nondelivery. It uses no
surrogate and no warm-start results.

\textbf{Recommendation rule.} For every baseline, if any online evaluation has
delivered by checkpoint $q$,
\begin{equation}
\rho_q = \arg\max_{k \le q:\; d(x_k) = 1} y(x_k),
\label{eq:recommendation}
\end{equation}
which is the best delivered online design. Otherwise the recommendation is an
untested design chosen by the method's own rule, the unqueried archive design
with the highest $s$ for the no-feedback method, with the highest selector
value for the selection-only method, and with the highest $m$ for the two
feedback-guided methods, and the next proposal for SMAC3. Ties break by a
stable candidate identifier. The checkpoint-32 recommendation is the terminal
one.

\subsection{Statistical procedures}

\textbf{Seed-block contrasts for the baselines.} For two methods $A$ and $B$
and seed $j$, the delivery effect at checkpoint 32 is the mean over the three
tasks $t$ of $d_{A,t,j} - d_{B,t,j}$, and the quality effect is the mean of
$y_{A,t,j} - y_{B,t,j}$ over the tasks where both recommendations deliver,
undefined when there is none. With $n_b$ available blocks and effects
$\Delta_1, \dots, \Delta_{n_b}$, the estimate is $\bar{\Delta}$ and the exact
two-sided sign-flip $p$ value is
\begin{equation}
p = 2^{-n_b} \,\bigl| \{ \varepsilon \in \{-1, +1\}^{n_b} :
\bigl| \tfrac{1}{n_b} \textstyle\sum_j \varepsilon_j \Delta_j \bigr| \ge
| \bar{\Delta} | \} \bigr| .
\label{eq:signflip}
\end{equation}
The 95\% interval is the 2.5th and 97.5th percentile of $\bar{\Delta}$ over
10{,}000 resamples of the blocks with replacement, drawn from a PCG64
generator with seed 73001, fixed in the analysis plan before any result was
known. The six planned tests are corrected by Holm's step-down procedure. With
ordered $p$ values $p_{(1)} \le \dots \le p_{(6)}$ the adjusted values are
$\tilde{p}_{(i)} = \max_{j \le i} \min\{1, (7 - j)\, p_{(j)}\}$, and a
contrast is significant when $\tilde{p} \le 0.05$. With seven blocks the
smallest attainable $p$ in Equation~\ref{eq:signflip} is $2/128 = 0.016$,
which is why Section~\ref{sec:results} reports the quality contrasts as
underpowered, not null.

\textbf{Agent contrasts.} A seed label does not pair agent episodes, so the two
methods' ten episodes are resampled independently, 10{,}000 times with the same
generator. The quality difference is the difference between the two
conditional means of Equation~\ref{eq:endpoints}, and each resample
recomputes both means from its delivered recommendations.
Table~\ref{tab:agent-contrasts} lists them.

\textbf{Intervals for a single rate.} For $k$ delivered recommendations out of
$n$, Table~\ref{tab:marginal} reports the exact Clopper--Pearson interval
\begin{equation}
\bigl[ B_{\alpha/2}(k,\, n - k + 1),\; B_{1 - \alpha/2}(k + 1,\, n - k) \bigr],
\label{eq:clopper}
\end{equation}
where $B_{u}(a, b)$ is the $u$ quantile of the beta distribution, the ends
are 0 when $k = 0$ and 1 when $k = n$, and $\alpha = 0.05$. Its coverage is at
least the nominal 95\% for every $k$.

\textbf{Intervals for a difference of rates.} A bootstrap of two arms that both
deliver every recommendation has no spread, so Table~\ref{tab:agent-contrasts}
uses Newcombe's score interval. With $\hat{p} = k / n$, $z = 1.96$, and the
Wilson bounds
\begin{equation}
l, u = \frac{\hat{p} + z^{2}/2n \;\mp\; z \sqrt{\hat{p}(1 - \hat{p})/n + z^{2}/4n^{2}}}
{1 + z^{2}/n},
\label{eq:wilson}
\end{equation}
the interval for $\hat{p}_1 - \hat{p}_2$ is
\begin{equation}
\Bigl[ \hat{p}_1 - \hat{p}_2 - \sqrt{(\hat{p}_1 - l_1)^{2} + (u_2 - \hat{p}_2)^{2}},\;
\hat{p}_1 - \hat{p}_2 + \sqrt{(u_1 - \hat{p}_1)^{2} + (\hat{p}_2 - l_2)^{2}} \Bigr].
\label{eq:newcombe}
\end{equation}

\subsection{Two remarks}

\textbf{Why the surrogate is optimistic where the search goes.} Write the
surrogate score of candidate $i$ as $s_i = y_i + \epsilon_i$ with an error
$\epsilon_i$ that has mean zero for a randomly chosen design. A search that
screens many candidates and keeps $i^{*} = \arg\max_i s_i$ selects for large
positive errors as much as for large true scores. This is the optimizer's
curse of \citet{smith2006curse}. Under their assumption that the error has mean
zero given the true score, selecting the maximum gives
$\mathbb{E}[\, s_{i^{*}} - y_{i^{*}} \,] \ge 0$, and the gap grows with the
number of candidates screened and with the error variance. In WingBench the error variance is
largest away from the SuperWing data, which is where optimization pushes, and
the same selection applies to delivery, since a design whose failure the
surrogate does not see is favored exactly because that failure is not priced.
Appendix~\ref{app:classifier} measures the consequence, and the agents that
fitted their own models in Section~\ref{sec:results} met the same effect.

\textbf{Why the recommendation rule is the best delivered design.} Repeated
evaluations of the same wing agreed in 800 of 809 cases
(Appendix~\ref{app:outcomes}), so a private regrade of an already delivered
design returns its online score. Among the designs an optimizer has evaluated,
Equation~\ref{eq:recommendation} therefore maximizes the delivered score with
certainty. An untested design $x$ can only beat it in expectation when
$p(x)\, \mathbb{E}[\, y(x) \mid d(x) = 1 \,]$ exceeds the best delivered score,
and the baselines never take that gamble once a design has delivered. This is
a conservative rule, chosen because the benchmark asks for a design worth
handing to a user, not for the design with the best expected score.

\begin{table}[H]
\caption{Settings of the reference methods and of the analysis, as fixed in
the released configurations.}
\label{tab:app-settings}
\centering
\scriptsize
\begin{tabular*}{\textwidth}{@{\extracolsep{\fill}}ll}
\toprule
setting & value \\
\midrule
design box & $[-3.5, 3.5]^{38}$ standardized, edge tolerance $10^{-9}$ \\
flow conditions & Mach $0.7853$, $0.8195$, $0.8530$ at $\mathrm{Re} = 2 \times 10^{7}$ \\
lift match & $|C_L - C_L^{*}| \le 0.02$, lift oscillation $< 5 \times 10^{-4}$, $\alpha \in [-10^{\circ}, 25^{\circ}]$ \\
score envelope & $0 < r_c \le 40$ for every condition \\
time limits & 30 min per mesh, 2 h per flow condition, wall clock \\
reference machine & one job at a time on an AWS m7a.4xlarge, 16 cores, 64 GB \\
surrogate budget & 4{,}096 evaluations in batches of at most 16 \\
online budget & 32 sequential CFD evaluations \\
checkpoints & 0, 1, 2, 4, 8, 16, 32 \\
warm start & the same 62 CFD results for the four surrogate-assisted baselines \\
population, generations & 16, 256 (eight generations per online query when interleaved) \\
active archive & 64 designs \\
CMA-ES & initial mean $0.5$ in unit coordinates, $\sigma_0 = 0.25$, no mirrors \\
DE & DE/rand/1/bin, $F = 0.8$, $\mathit{CR} = 0.9$ \\
invalid surrogate value & loss $10^{12}$ \\
delivery model & logistic, $\ell_2$ penalty 4 (intercept 0.1), 40 iterations, clip to $[0.05, 0.95]$ \\
outcome weights & pass 1, conclusive failure 1, partial 0.5, undetermined 0.5 \\
discrepancy model & ridge in $\log y - \log s$, penalty 10 (intercept 0.1), $\hat{\sigma} \ge 0.05$, 0.25 with no label \\
merit & $\log p + \log s + \hat{\delta} + 0.5\, \hat{\sigma} + 0.1\, H(p)$ \\
SMAC3 & version 2.4.0, black-box facade, 8 Sobol points, 24 model-guided queries \\
agents & DeepSeek V4.1 Flash, maximum reasoning effort, 64 tool rounds per turn \\
reflection & summary capped at 4{,}000 characters \\
planning & four plans of two steps, at most three cycles per decision \\
bootstrap & 10{,}000 resamples, PCG64 seed 73001, 95\% percentile intervals \\
multiple comparisons & Holm over the six planned baseline contrasts at $\alpha = 0.05$ \\
\bottomrule
\end{tabular*}
\end{table}
